\section{Introduction}\label{sec:intro}

AI agents increasingly interact with external systems through tools. The Model Context Protocol (MCP) allows assistants to read email, modify files, query databases, and execute commands using tools provided by third-party servers. Users authorize these interactions based on the information exposed by the server, primarily the tool's name and natural-language description. This rests on a trust assumption: the capabilities available to the agent remain consistent with those presented at approval.

However, MCP servers continue to change after approval, and those changes do not necessarily trigger renewed authorization. In September 2025, an MCP email server was reported to have released fifteen benign versions before introducing an update that silently copied outgoing messages to an address controlled by its author~\cite{koi2025postmark}. Previously approved users received no new approval request. This problem extends beyond individual implementations. Rashidi~\cite{unicodetag2026} found that none of eight tool-metadata attack techniques, including changes introduced after an initially benign approval, triggered mandatory re-approval under the protocol.

These findings establish the possibility of exploiting persistent approval, but they leave a broader empirical question unanswered: \textbf{how do MCP tools actually evolve after approval, and how often do their capabilities change without corresponding changes to their descriptions?} Existing large-scale studies of MCP security examine server behavior at a single point in time~\cite{dcichecker2026,mcpdiff2026,mcpfsd2026,mcpzoo2026}. Such snapshots reveal what a server can do at a particular version, but cannot establish when capabilities were introduced, whether their descriptions changed, or whether existing approvals would remain valid.

To address this gap, we conduct a longitudinal analysis of 2,571 open-source MCP servers across 35,725 sampled versions. For each version, we extract tool descriptions and identify code-reachable capabilities, including shell execution, dynamic code evaluation, file reads and writes, secret access, and network operations. We then compare successive versions to identify changes in tool behavior and whether those changes are reflected in the descriptions exposed to users. We validate our measurements through manual annotation of a stratified sample and replay observed release histories under alternative client re-approval policies.

Our main findings are:
\begin{itemize}
\item \textbf{Approved servers change continually.} 65.7\% of servers edit the description of a tool that already existed, and 19.5\% give an existing tool a capability class it did not have before, typically weeks after the first release (Section~\ref{sec:rq1}).
\item \textbf{Most of that change is invisible.} 73.9\% of capability gains leave the tool's description byte-identical, and only 2.1\% of high-risk gains are mentioned in the description users approved. In 3.5\% of servers a high-risk capability arrives with no visible change at all (Section~\ref{sec:rq2}).
\item \textbf{Most changes are benign, which is the problem.} In a hand-labeled sample, about half of the flagged changes are real, and 30\% of the real ones are security-relevant: new access to credentials, AI-controlled execution or file writes, new local channels. Risky changes are rare events inside constant routine updates (Section~\ref{sec:rq3}).
\item \textbf{The attacks studied so far are not what happens.} Hidden instructions in descriptions are virtually absent (one server in 2,571). Instead, a few servers move risky actions into the agent itself, by instructing the model to run commands or by imitating the client's own system messages, where no server-side scan looks (Section~\ref{sec:cases}).
\item \textbf{Description-based re-approval fails.} Re-approving a tool when its description changes surfaces none of the silent capability gains; re-approving on every update prompts a median of 15.7 times per server per year. Rules tied to code-derived capabilities prompt least and target the changes that matter (Section~\ref{sec:rq4}).
\end{itemize}

In summary, we make the following contributions:
\begin{itemize}
\item the first longitudinal dataset of MCP servers, covering 35,725 sampled versions of 2,571 servers, with per-version tools, descriptions, and code-derived capabilities;
\item a measurement of how often, how visibly, and how riskily approved servers change, validated against human annotation;
\item a characterization of agent-mediated execution, a drift pattern invisible to server-side analysis;
\item a policy analysis on real histories, and concrete recommendations for MCP clients and the protocol.
\end{itemize}
We will release our code and data (Appendix~\ref{sec:appendix}).
